% Direct-authority diplomatic French for printed p.207.
\clearpage
\SourcePage{207}

\LayerTarget{tohoku:remark:smith-fixed-point-application:p1}
\noindent dans [2] est entièrement inutile. On trouve encore plus facilement,
comme application du corollaire précédent, que si un groupe fini $G$ d'ordre
premier $p$ opère dans un espace $X$ de dimension finie qui est acyclique mod.
$p$, alors l'ensemble $X_0$ des points fixes et l'espace quotient $X/G$ sont
acycliques mod. $p$. Il en résulte que si $X$ est acyclique pour un corps de
coefficients $k$ (de caractéristique quelconque, le cas de caractéristique
$\ne p$ étant trivial), il en est de même de $X/G$ ; et $k$ étant arbitraire,
cet énoncé reste valable en y remplaçant $k$ par $\mathbf Z$. Une récurrence
immédiate montre que ces résultats relatifs à $X/G$ sont encore valables si on
suppose seulement $G$ fini résoluble. Enfin, les calculs de ce numéro sont
aussi utiles dans l'étude des puissances de Steenrod dans les faisceaux et la
cohomologie des puissances symétriques d'un espace. Ce sont d'ailleurs des
calculs analogues dans un manuscrit de R. Godement sur les puissances de
Steenrod (non publié), qui ont inspiré le présent numéro.

\LayerTarget{tohoku:section:5.5}
\bigskip
\noindent\textbf{5.5. Calcul des $H^n(X;G,A)$ par les recouvrements.} Nous
aurons à nous servir de résultats auxiliaires :

\LayerTarget{tohoku:lemma:5.5.1}
\noindent\textsc{Lemme 5.5.1.} \textit{Soient $C_1,C_2,C_3$ trois catégories
abéliennes, on suppose que dans les deux premières tout objet est isomorphe à
un sous-truc d'un objet injectif. Soient $F$ et $G$ des foncteurs covariants de
$C_1$ dans $C_2$ et de $C_2$ dans $C_3$ respectivement, on suppose que $F$
transforme objets injectifs en objets $G$-acycliques. Soit enfin $\mathbf K$ un
foncteur covariant de $C_1$ dans la catégorie $\mathbf K(C_2)$ des complexes à
degrés positifs dans $C_2$, satisfaisant aux deux conditions (i)
$H^0\mathbf K(A)=F(A)$ (isomorphisme fonctoriel) (ii) $\mathbf K(A)$ est
acyclique en dimensions $n>0$ quand $A$ est injectif. Sous ces conditions, il
existe sur $C_1$ un foncteur spectral cohomologique, aboutissant au foncteur
gradué $(R^n(GF)(A))$, et dont le terme initial est}

\LayerTarget{tohoku:equation:5.5.1}
\[
\tag{5.5.1}
E_2^{p,q}(A)=R^pG\bigl(R^q\mathbf K(A)\bigr).
\]

\LayerTarget{tohoku:proof:lemma-5.5.1}
Désignons par $H^0$ le foncteur covariant $L\longmapsto H^0(L)$ de
$\mathbf K(C_2)$ dans $C_2$, on a alors
$GF=G(H^0\mathbf K)=(GH^0)\mathbf K$. Or si $A$ est injectif, $\mathbf K(A)$
est $(GH^0)$-acyclique, i.e. on a $R^nG(\mathbf K(A))=0$ pour $n>0$ ; en
effet, $\mathbf K(A)$ est une résolution de $F(A)$ en vertu de (i) et (ii),
donc $R^nG(\mathbf K(A))=R^nG(F(A))$ et le dernier membre est nul en vertu de
l'hypothèse faite sur $F$. On peut donc appliquer le théorème 2.4.1 au
foncteur composé $(GH^0)\mathbf K$, ce qui donne le foncteur spectral du lemme.

\LayerTarget{tohoku:lemma:5.5.2}
\noindent\textsc{Lemme 5.5.2.} \textit{$C_1,C_2,C_3,F$ et $G$ étant comme
dans le lemme précédent, il existe un foncteur spectral cohomologique sur la
catégorie $\mathbf K(C_1)$ des complexes $C$ à degrés positifs dans $C_1$,
aboutissant au foncteur gradué $(R^n(GF)(C))$, et dont le terme initial est}

\LayerTarget{tohoku:equation:5.5.2}
\[
\tag{5.5.2}
E_2^{p,q}(C)=R^pG\bigl(R^qF(C)\bigr).
\]

\LayerTarget{tohoku:proof:lemma-5.5.2:p1}
Indiquons par une barre au dessus d'un foncteur, le foncteur obtenu par
``prolongement'' aux catégories de complexes à degrés positifs. On a alors
\[
R^n(GF)=R^n\bigl(H^0\overline{GF}\bigr)
=R^n\bigl((H^0\overline G)\overline F\bigr).
\]
Or le foncteur $\overline F:\mathbf K(C_1)\longrightarrow\mathbf K(C_2)$
transforme objets injectifs en objets $(H^0\overline G)$-acycliques, en
d'autres termes, si $C$ est un objet injectif de $\mathbf K(C_1)$, alors
$R^nG(\overline F(C))=0$ pour $n>0$ ; en effet, d'après la cara-

\TohokuContinuation{208}
